\documentclass[10pt,a4paper]{amsart}
\usepackage[margin=19mm]{geometry}
\usepackage{amsmath,amssymb,mathtools}
\usepackage[scr=rsfs,cal=euler]{mathalfa}
\usepackage{xfrac}
\usepackage[unicode,hidelinks,bookmarks=false]{hyperref}
\newcommand{\ie}{i.e.\ }
\newcommand{\cf}{cf.\ }
\newcommand{\resp}{respectively\ }
\newcommand{\Subsection}{\subsection}
\providecommand{\nobreakdash}{\nobreak\hbox{-}}
\setlength{\emergencystretch}{2em}
\multlinegap0pt
\usepackage{siunitx}
\usepackage{siunitx}
\usepackage{siunitx}
\usepackage{siunitx}
\usepackage{tikz}
\usepackage{mathtools}
\usepackage{cancel}
\usepackage{xcolor}
\usepackage{siunitx}
\usepackage{amssymb}
\usepackage{float}
\usepackage[shortlabels]{enumitem}
\newtheorem{theorem}{Theorem}
\usepackage{cleveref}
\crefname{theorem}{Theorem}{Theorems}
\DeclareMathOperator{\BR}{BR}
\newcommand{\R}{\mathbb{R}}
\newcommand{\cC}{\mathcal{C}}
\newcommand{\lilc}{locally inner Lipschitz continuous}
\title{On the Connectedness of Sublevel Sets in Invex Optimization}
\author{Vinzenz Thoma, Zebang Shen, Niao He}
\date{Source: arXiv:2604.12045v1 (CC BY 4.0)}
\begin{document}
\maketitle

\noindent\emph{Source and licence.} On the Connectedness of Sublevel Sets in Invex Optimization. Version arXiv:2604.12045v1 (CC BY 4.0), distributed by arXiv under the Creative Commons Attribution 4.0 International licence, http://creativecommons.org/licenses/by/4.0/, as recorded in the arXiv licence metadata for this exact version and shown on the abstract page https://arxiv.org/abs/2604.12045v1. Full licence notice: \texttt{licenses/arxiv-2604.12045-CC-BY-4.0.txt}.\par\smallskip

\noindent\textit{Editorial scope.} Counted unit: the article's theorem thm:Holderconnected (source0.tex lines 1223--1226). The article's complete original proof of it is delivered with it (source0.tex lines 1229--1272, one \texttt{proof} environment of 44 lines). No mathematics is written, repaired or re-derived: the statement and the proof are the article's own lines, in the article's own notation, and no retained line is substituted or re-flowed.  Retained uncounted as the source-local context the proof needs: lines 423--425.  Nothing was omitted from any retained span, so the package carries no omission record.  No retained line contains a citation, so the wrapper carries no bibliography and no bibliography entry.  No retained line contains a figure environment, an image-inclusion command or drawing code, so no image is delivered and the package declares no asset dependency.  Editorial formatting only, in addition: an AMS \texttt{amsart} preamble that restates the article's own declarations and macros used by the retained lines, namely the environments \texttt{theorem} on separate counters, together with the standard packages those lines need.  The retained environments print in this document as Theorem 1 in order of appearance, whereas the article prints the same items with the numbering of its own full document.  The complete native source of the article as distributed in the arXiv e-print for this exact version is bundled unchanged as \texttt{sources/198379/source0.tex}.
\begin{theorem}
\label{thm:lilc} Assume that for all $(x,y)\in \underline{M}$, $\BR_y(x)$ is \lilc\ at $x$ for $y$ and conversely for all $(x,y)\in \overline{M}$, $\BR_x(y)$ is \lilc\ at $y$ for $x$. Then $E=\underline{M}=\overline{M}$ and this common set is connected.
\end{theorem}

\begin{theorem}
\label{thm:Holderconnected}
Assume $f$ is continuously differentiable and that $\nabla f$ is locally Lipschitz continuous on $\underline{M}\cup\overline{M}$. Furthermore, assume that for all $(\overline{x},\overline{y})\in \underline{M}$, $\BR_y(x)$ is locally inner $\alpha_x$-Hölder continuous at $\overline{x}$ for $\overline{y}$ with $\alpha_x>0.5$ and conversely for all $(\overline{x},\overline{y})\in \overline{M}$, $\BR_x(y)$ is locally inner $\alpha_y$-Hölder continuous at $\overline{y}$ for $\overline{x}$ with $\alpha_y>0.5$. Then $E=\underline{M}=\overline{M}$ and this common set is connected.
\end{theorem}

\begin{proof}
    Analogous to the proof of \Cref{thm:lilc}, we restrict to showing that the local inner Hölder continuity of $\BR_y(x)$ implies that $\underline{M} \subseteq E$. The proof for $\overline{M}$ and thus the connectedness of $E$ follows analogously.

    Let $\overline{x}\in X$ and choose any $\overline{y}\in \BR_y(\overline{x})$. By the local inner Hölder continuity of $\BR_y(x)$, there exists a constant $\kappa_x$ and an open neighborhood $N_{\overline{x}}$ on which $\BR_y(x)$ is inner Hölder continuous. Let $\delta'$ be small enough, such that for all $\delta\leq \delta': \mathbb{B}_\delta (\overline{x})\subseteq N_{\overline{x}}$, where $\mathbb{B}_\delta (\overline{x})$ is an open ball of radius $\delta$ around $\overline{x}$, and let $v \in \R^{n_x}$ be a vector such that $\|v\|=1$. We choose a $y(\delta,v) \in \BR_y(\overline{x}+\delta v)$, such that $\|y(\delta,v)-\overline{y}\|=d(\overline{y},\BR_y(\overline{x}+\delta v))$. The latter is possible as $\BR_y(\overline{x}+\delta v)$ is compact.

    By the local inner Hölder continuity of $\BR_y({x})$ at $\overline{x}$ for $\overline{y}$, it holds that:
    \begin{align*}
       \|y(\delta,v)-\overline{y}\| =d(\overline{y},\BR_y(\overline{x}+\delta v)) \leq \kappa_x \|\delta v\|^{\alpha_x} = \kappa_x \delta^{\alpha_x}.
    \end{align*}

    Let $N^L$ be the neighborhood of $(\overline{x},\overline{y})$ on which $\nabla f$ is Lipschitz continuous, then if necessary we can choose a smaller $\delta'$, to ensure that for all $\delta\leq \delta': \mathbb{B}_\delta (\overline{x})\times \mathbb{B}_\delta (\overline{y})\subseteq N^L$. 
    

    Using that $f$ is $\cC^1$, we can derive the following expression for the change of $F$ along $v$:
    \begin{align*}
        F(\overline{x}+\delta v) - F(x) =& f(\overline{x} +\delta v,y(\delta,v)) - f(\overline{x},\overline{y}) \\
        =& \nabla_x f(\overline{x},\overline{y})^T \delta v + \underbrace{\nabla_y f(\overline{x},\overline{y})^T}_{=0} (y(\delta,v)-\overline{y}) + R_1(\delta,v)\\
        =& \nabla_x f(\overline{x},\overline{y})^T \delta v + R_1(\delta,v).
    \end{align*}

    By the fundamental theorem of calculus, the remainder term $R_1(\delta,v)$ is given by:
    \begin{align*}
        R_1(\delta,v) =& \int_0^1  \left(\nabla f(\overline{x}+t \delta v,\overline{y}+t(y(\delta,v)-\overline{y}))-\nabla f(\overline{x},\overline{y})\right)^T \left( \delta v, y(\delta,v)-\overline{y}\right) dt 
    \end{align*}

    Using the Lipschitz continuity of $\nabla f$ on $N^L$ and denoting the Lipschitz constant as $S$, we can bound the remainder term:
    \begin{align*}
        \|R_1(\delta,v)\| \leq & \int_0^1 \|\nabla f(\overline{x}+t \delta v,\overline{y}+t(y(\delta,v)-\overline{y}))-\nabla f(\overline{x},\overline{y})\|  \| \left( \delta v, y(\delta,v)-\overline{y}\right) \| dt \\
        \leq & S \int_0^1 t \|\left( \delta v, y(\delta,v)-\overline{y}\right) \|^2 dt \\
        = & \frac{S}{2} \|\left( \delta v, y(\delta,v)-\overline{y}\right) \|^2 \\
        = & \mathcal{O}(\delta^2 + \delta^{2\alpha_x} )\\
    \end{align*}
    Therefore, independent of the choice of $\overline{y}\in \BR_y(\overline{x})$, we have:
    \begin{align*}
        \lim_{\delta \to 0} \frac{F(\overline{x}+\delta v)-F(\overline{x})}{\delta} =& \nabla_x f(\overline{x},\overline{y})^T v + \lim_{\delta \to 0} \mathcal{O}(\delta+\delta^{2\alpha_x-1})\\
    \end{align*}    
    By Assumption $\alpha_x>0.5$, which implies that the extra term vanishes as $\delta \to 0$. Note, that this result holds for all $\overline{y}\in \BR_y(\overline{x})$. Therefore, all directional derivatives exist, independent of the choice of $\overline{y}\in \BR_y(\overline{x})$, and are continuous functions because $f$ is $\cC^1$. $F$ is therefore differentiable at $\overline{x}$ and its gradient is given by:
    \[
    \nabla F(\overline{x}) = \nabla_x f(\overline{x},\overline{y}),
    \]
    for any $\overline{y} \in \BR_y(\overline{x})$.\footnote{Indeed, from our analysis it follows that $\forall y \in \BR_y(\overline{x})\; \forall \overline{y} \in \BR_y(\overline{x}):\nabla_x f(\overline{x},y)=\nabla_x f(\overline{x},\overline{y})$.} 

    To conclude our proof, consider an arbitrary point $(\overline{x},\overline{y})\in \underline{M}$. By definition, $\nabla_y f(\overline{x},\overline{y})=0$. As $\overline{x}\in \arg\min_x F(x)$, it follows that $\nabla F(\overline{x})=0$ and therefore, using our result above, $\nabla_x f(\overline{x},\overline{y})=\nabla F(\overline{x})=0$. Due to the invex-incave structure of $f$, this implies $(\overline{x},\overline{y})\in E$ and thus $\underline{M}\subseteq E$.
\end{proof}

\end{document}
